\begin{figure}[htbp]
\centering
\begin{tikzpicture}
\begin{axis}[
    width=12cm, height=8cm,
    xlabel={$x$},
    ylabel={$f(x)$},
    xmin=-0.5, xmax=4.5,
    ymin=-0.5, ymax=5.5,
    grid=major,
    grid style={UofSC50Black!30},
    axis lines=middle,
    xtick={1,2,3,4},
    ytick={1,2,3,4,5},
    clip=false
]

% The function f(x) = 0.3x^2 + 0.5
\addplot[UofSCBlack, very thick, domain=0:4.2, samples=100] {0.3*x^2 + 0.5};
\addlegendentry{$f(x)$}

% Point x_0 at x=1.5
\pgfmathsetmacro{\xzero}{1.5}
\pgfmathsetmacro{\yzero}{0.3*\xzero*\xzero + 0.5}
\pgfmathsetmacro{\slope}{0.6*\xzero}

% Point x_0 + Delta x at x=3.0
\pgfmathsetmacro{\xone}{3.0}
\pgfmathsetmacro{\yone}{0.3*\xone*\xone + 0.5}

% Secant line between the two points
\pgfmathsetmacro{\secantslope}{(\yone-\yzero)/(\xone-\xzero)}
\addplot[UofSCAtlantic, thick, dashed, domain=0.5:4, samples=2] {\yzero + \secantslope*(x-\xzero)};
\addlegendentry{Secant line}

% Tangent line at x_0 (derivative)
\addplot[UofSCGarnet, very thick, domain=0:3.5, samples=2] {\yzero + \slope*(x-\xzero)};
\addlegendentry{Tangent line}

% Points on curve
\node[circle, fill=UofSCGarnet, minimum size=6pt, inner sep=0pt] at (axis cs:\xzero,\yzero) {};
\node[circle, fill=UofSCAtlantic, minimum size=6pt, inner sep=0pt] at (axis cs:\xone,\yone) {};

% Labels for points
\node[below left, UofSCGarnet, font=\small\bfseries] at (axis cs:\xzero,\yzero) {$(x_0, f(x_0))$};
\node[above right, UofSCAtlantic, font=\small\bfseries] at (axis cs:\xone,\yone) {$(x_0 + \Delta x, f(x_0 + \Delta x))$};

% Delta x bracket
\draw[UofSCHorseshoe, thick, |<->|] (axis cs:\xzero,0.2) -- (axis cs:\xone,0.2);
\node[below, UofSCHorseshoe, font=\bfseries] at (axis cs:{(\xzero+\xone)/2},0.2) {$\Delta x$};

% Delta f bracket (vertical)
\draw[UofSCCongaree, thick, |<->|] (axis cs:3.3,\yzero) -- (axis cs:3.3,\yone);
\node[right, UofSCCongaree, font=\bfseries] at (axis cs:3.3,{(\yzero+\yone)/2}) {$\Delta f$};

% Horizontal and vertical dashed lines for visualization
\draw[UofSC50Black, thin, dashed] (axis cs:\xzero,\yzero) -- (axis cs:\xone,\yzero);
\draw[UofSC50Black, thin, dashed] (axis cs:\xone,\yzero) -- (axis cs:\xone,\yone);

\end{axis}

% Derivative formula annotation
\node[anchor=north west, align=left, font=\small] at (0.5,1.5) {
    \textcolor{UofSCAtlantic}{Secant slope:} $\displaystyle\frac{\Delta f}{\Delta x} = \frac{f(x_0 + \Delta x) - f(x_0)}{\Delta x}$\\[8pt]
    \textcolor{UofSCGarnet}{Derivative:} $\displaystyle f'(x_0) = \lim_{\Delta x \to 0} \frac{\Delta f}{\Delta x}$
};

\end{tikzpicture}
\caption{Visualization of the derivative as the limit of secant line slopes. The curve $f(x)$ is shown in black. The secant line (dashed Atlantic blue) connects two points separated by $\Delta x$, with slope $\Delta f/\Delta x$. As $\Delta x \to 0$, the secant line approaches the tangent line (Garnet red) at $x_0$, whose slope equals the derivative $f'(x_0)$. The intervals $\Delta x$ (Horseshoe green) and $\Delta f$ (Congaree teal) are marked.}
\label{fig:derivative_visualization}
\end{figure}
